\chapter{Interface e Experiência do Usuário (UI/UX) Científica}
\label{sec:uiux}

\section{Design da Interface e Fluxo de Trabalho}

A interface gráfica do sistema modernizado foi desenvolvida sobre o \textit{framework} Streamlit, adotando princípios de usabilidade voltados para aplicações científicas e de análise de dados. A organização visual busca minimizar a poluição de tela e guiar o usuário de forma intuitiva pelas etapas do processo de calibração multivariada.

A aplicação é dividida em três abas sequenciais que espelham o fluxo de trabalho do usuário:
\begin{enumerate}[label=\alph*)]
    \item \textbf{Aba 1 (Configuração e Carga de Dados):} Destinada ao envio dos arquivos de calibração e teste nos formatos \texttt{.txt}, \texttt{.csv} ou \texttt{.xlsx}, e à definição dos hiperparâmetros de entrada (número de replicações \textit{m} e fração de subamostragem \textit{k});
    \item \textbf{Aba 2 (Resultados Numéricos e Logs):} Exibe as métricas de avaliação preditiva (RMSEP, MAE, $R^2$), a indicação dos parâmetros ótimos selecionados pela validação cruzada 5-\textit{fold} e o \textit{log} detalhado de execução;
    \item \textbf{Aba 3 (Visualização Gráfica):} Apresenta os gráficos interativos gerados via Plotly, incluindo os gráficos de paridade (Valores Reais \textit{versus} Preditos) e as superfícies de erro.
\end{enumerate}

Essa estrutura em abas sequenciais organiza o \textit{workflow} analítico segundo as etapas operacionais do usuário, reduzindo a dispersão visual e delimitando os passos de execução.

O ambiente de implantação viabiliza o acesso simultâneo via navegador sem necessidade de instalação local de dependências ou ambientes de execução Java, sujeitando-se aos limites de concorrência e memória da plataforma de hospedagem adotada.